% =====================================================================
%  preamble.tex  --  SCGP 四维流形讲义（中文精译）通用导言区
%  编译方式： XeLaTeX（推荐）或 pdfLaTeX + ctex 宏包
% =====================================================================

% ---------- 页面与版式 ----------
\usepackage[a4paper, top=2.6cm, bottom=2.6cm, left=2.7cm, right=2.7cm]{geometry}
\usepackage[UTF8, fontset=fandol, scheme=plain]{ctex}
\setlength{\parskip}{0.32em}
\linespread{1.22}

% ---------- 数学 ----------
\usepackage{amsmath, amssymb, amsthm, mathtools}
\usepackage{mathrsfs}
\usepackage{bm}
\usepackage{dsfont}

% ---------- 图形与表格 ----------
\usepackage{graphicx}
\usepackage{float}
\usepackage{array}
\usepackage{booktabs}
\usepackage{caption}
\captionsetup{font=small, labelsep=quad}
\renewcommand{\figurename}{图}
\renewcommand{\tablename}{表}
\renewcommand{\contentsname}{目录}
\renewcommand{\refname}{参考文献}
\renewcommand{\abstractname}{摘要}
\renewcommand{\appendixname}{附录}

% ---------- 颜色 / 超链接 / 列表 ----------
\usepackage{xcolor}
\usepackage[colorlinks=true, linkcolor=blue!45!black, citecolor=green!40!black,
            urlcolor=magenta!55!black, bookmarksnumbered=true,
            hypertexnames=false]{hyperref}
\usepackage{enumitem}
\setlist{itemsep=2pt, topsep=3pt, parsep=0pt}

% ---------- 页眉页脚 ----------
\usepackage{fancyhdr}
\pagestyle{fancy}
\fancyhf{}
\fancyhead[LE,RO]{\small\thepage}
\fancyhead[LO]{\small\nouppercase{\rightmark}}
\fancyhead[RE]{\small\nouppercase{\leftmark}}
\renewcommand{\headrulewidth}{0.3pt}

% ---------- 定理类环境 ----------
\theoremstyle{plain}
\newtheorem{theorem}{定理}[section]
\newtheorem{proposition}[theorem]{命题}
\newtheorem{lemma}[theorem]{引理}
\newtheorem{corollary}[theorem]{推论}
\newtheorem{conjecture}[theorem]{猜想}
\theoremstyle{definition}
\newtheorem{definition}[theorem]{定义}
\newtheorem{example}[theorem]{例}
\newtheorem{remark}[theorem]{注}
\theoremstyle{remark}
\newtheorem*{remark*}{注}

% ---------- 常用记号宏（全书统一） ----------
\newcommand{\Tr}{\operatorname{Tr}}
\newcommand{\tr}{\operatorname{tr}}
\newcommand{\Aut}{\operatorname{Aut}}
\newcommand{\Conn}{\operatorname{Conn}}
\newcommand{\Sym}{\operatorname{Sym}}
\newcommand{\Tors}{\operatorname{Tors}}
\newcommand{\Coker}{\operatorname{Cok}}
\newcommand{\Ind}{\operatorname{Ind}}
\newcommand{\vdim}{\operatorname{vdim}}
\newcommand{\sign}{\operatorname{sign}}
\newcommand{\Stab}{\operatorname{Stab}}
\newcommand{\Lie}{\operatorname{Lie}}
\newcommand{\Ber}{\operatorname{Ber}}
\newcommand{\diag}{\operatorname{diag}}
\newcommand{\ad}{\operatorname{ad}}
\newcommand{\im}{\operatorname{Im}}
\newcommand{\coker}{\operatorname{Cok}}
\newcommand{\rk}{\operatorname{rank}}
\newcommand{\wh}{\widehat}

\newcommand{\ZZ}{\mathbb{Z}}
\newcommand{\RR}{\mathbb{R}}
\newcommand{\CC}{\mathbb{C}}
\newcommand{\HH}{\mathbb{H}}
\newcommand{\PP}{\mathbb{P}}
\newcommand{\EE}{\mathbb{E}}
\newcommand{\KK}{\mathbb{K}}
\newcommand{\OO}{\mathcal{O}}
\newcommand{\cO}{\mathcal{O}}
\newcommand{\cA}{\mathcal{A}}
\newcommand{\cE}{\mathcal{E}}
\newcommand{\cM}{\mathcal{M}}
\newcommand{\cN}{\mathcal{N}}
\newcommand{\cF}{\mathcal{F}}
\newcommand{\cC}{\mathcal{C}}
\newcommand{\fs}{\mathfrak}
\newcommand{\sfs}[1]{\mathfrak{#1}}
\newcommand{\su}{\mathfrak{su}}
\newcommand{\suf}{\mathfrak{u}}
\newcommand{\so}{\mathfrak{so}}
\newcommand{\sog}{\mathfrak{g}}
\newcommand{\sot}{\mathfrak{t}}
\newcommand{\soh}{\mathfrak{h}}

\newcommand{\dd}{\mathrm{d}}
\newcommand{\ii}{\mathrm{i}}
\newcommand{\ee}{\mathrm{e}}
\newcommand{\vol}{\mathrm{vol}}
\newcommand{\Rdr}{H^{2}_{\mathrm{DR}}}
\newcommand{\cxn}{{\wh{\nabla}}}
\newcommand{\dev}[1]{\left\langle #1 \right\rangle}
\newcommand{\Avg}[1]{\left\langle #1 \right\rangle}

% 便于保留原文公式编号：\taggedeq{编号}{...} 或直接 equation+tag
\newcommand{\eqtag}[1]{\tag{#1}}

% ---------- 「本节重点」环境（每节开头） ----------
\newenvironment{keytakeaways}%
  {\par\smallskip\noindent\begin{minipage}{\linewidth}%
   \color{black!85}\hrule\smallskip\noindent\textbf{本节重点（Key Takeaways）}\par\smallskip%
   \begin{itemize}[leftmargin=1.6em, itemsep=1.5pt, topsep=1pt, parsep=0pt]}%
  {\end{itemize}\smallskip\hrule\end{minipage}\par\smallskip}

% 图占位（原 PDF 中仅有图题而无实际图形时使用）
\newcommand{\figplaceholder}[2][4.2cm]{%
  \fbox{\parbox[c][#1][c]{0.72\textwidth}{\centering\small\itshape #2}}}
